% ---
% titre: 10G2 - Nombres réels
% theme: NO
% chapitre: Nombres réels
% sous_chapitre:
% niveau: [10e]
% type: PLANIF
% tags: [c-nombres-reels,c-ensembles-nombres,c-fractions,c-puissance,c-racine,c-notation-scientifique,c-probabilites,c-nombres-decimaux]
% source: latex
% ---

\header{Nombres réels - Objectifs}
 
\def\arraystretch{1.2}
\begin{tabularx}{\linewidth}{|>{\centering\arraybackslash}p{3cm}|>{\centering\arraybackslash}p{3cm}|X|}
\hline
\cellcolor{themeNO!30} Livre & \cellcolor{themeNO!30} Fiches & \cellcolor{themeNO!30} \hfill Aide-mémoire\hfill\null\\
\hline
\multirow{5}{*}{pp. 35-71}&\multirow{5}{*}{pp. 32-67}&Généralités: pp. 24-26\\
&&Nombres rationnels: pp. 41-45\\
&&Nombres réels: p. 46\\
&&Puissances: pp. 46-48\\
&&Racines: pp. 50-51\\
\hline
\end{tabularx}
\def\arraystretch{1}

\def\arraystretch{1.6}
\begin{tabularx}{\textwidth}{|>{\raggedright\arraybackslash}X|>{\raggedright\arraybackslash}X|}
\hline
\cellcolor{themeNO!30}Compétences visées & \cellcolor{themeNO!30}Utilité dans la vie de tous les jours \\
\hline
Mobiliser les savoirs, savoir-faire et savoir être pour résoudre des problèmes en lien avec les nombres réels.
&\vspace{-1em}
\begin{itemize}[leftmargin=*, itemsep=2pt, topsep=0pt, partopsep=0pt, parsep=0pt, label=$\triangleright$]
    \item Évaluer les prix, les dépenses, le salaire
    \item Comprendre des statistiques et des probabilités
\end{itemize} \\
\hline
 
Développer des stratégies de recherche afin de résoudre des problèmes.
&\vspace{-1em}
\begin{itemize}[leftmargin=*, itemsep=2pt, topsep=0pt, label=$\triangleright$]
    \item Découper un projet en sous-tâches
    \item Choisir l'ordre des tâches à effectuer
\end{itemize} \\
\hline
 
Présenter sa démarche de façon claire (explications, grandeur cherchée, unités, calculs, réponse).
&\vspace{-1em}
\begin{itemize}[leftmargin=*, itemsep=2pt, topsep=0pt, label=$\triangleright$]
    \item Collaborer avec d'autres personnes
    \item Expliciter sa compréhension d'un sujet
    \item Argumenter
\end{itemize} \\
\hline
 
Développer des stratégies d'apprentissage pour s'approprier puis mobiliser les savoirs et savoir-faire.
&\vspace{-1em}
\begin{itemize}[leftmargin=*, itemsep=2pt, topsep=0pt, label=$\triangleright$]
    \item Apprendre une nouvelle activité %(art, sport, \ldots)
    \item S'adapter à un nouvel environnement %(travail, voyage, \ldots)
    \item Se former en autonomie %(tutoriel, lecture, \ldots)
\end{itemize} \\
\hline
 
Vérifier sa solution ou au moins sa plausibilité.
&\vspace{-1em}
\begin{itemize}[leftmargin=*, itemsep=2pt, topsep=0pt, label=$\triangleright$]
    \item S'autocritiquer
    \item Évaluer la qualité de son travail
\end{itemize} \\
\hline
\end{tabularx}
\def\arraystretch{1}
 
 
\vfill
\begin{astuce}{Ce que je dois déjà savoir}
\begin{itemize}[leftmargin=*, itemsep=2pt, topsep=2pt, label=$\triangleright$]
    \item Effectuer les \textcolor{themeNO}{quatre opérations} avec des nombres relatifs \hfill (NO 10.1)
    \item Connaître les priorités des opérations (\textcolor{themeNO}{PEMDAS}) \hfill (NO 10.1)
    \item Calculer avec des \textcolor{themeNO}{puissances} (base et exposant) \hfill (NO 10.1)
    \item Déterminer le \textcolor{themeNO}{ppmc} et le \textcolor{themeNO}{pgdc} de deux nombres \hfill (NO 10.1)
\end{itemize}
\end{astuce}

\newpage
%---------------
\null
\vfill
\def\arraystretch{1.6}
\begin{tabularx}{\textwidth}{|>{\raggedright\arraybackslash}p{5.2cm}|>{\raggedright\arraybackslash}X|>{\centering\arraybackslash}p{0.8cm}|>{\centering\arraybackslash}p{0.8cm}|>{\centering\arraybackslash}p{0.8cm}|}
\hline
\rowcolor{themeNO!30}\multicolumn{2}{|l|}{Objectifs} &\multicolumn{3}{c|}{Auto-évaluation} \\
\rowcolor{themeNO!30}\multicolumn{2}{|l|}{} & \Sadey[1.5]& \Neutrey[1.5] & \Smiley[1.5]  \\\hline
 
\multirow{8}{=}{Connaître et utiliser les différentes écritures d'un même nombre. Comparer, approximer, encadrer, représenter sur une droite et déterminer l'ordre de grandeur de nombres écrits sous forme~:
\begin{itemize}[leftmargin=*, itemsep=1pt, topsep=2pt, label=$\triangleright$]
    \item fractionnaire (y compris simplification / amplification) dans $\mathbb{Q}$
    \item de pourcentage
    \item de puissance $a^b$ ($a$ sous forme décimale dans $\mathbb{Q}$~; $b$ dans $\mathbb{N}$)
    \item de notation scientifique $a\cdot 10^n$ ($n$ dans $\mathbb{Z}$)
    \item de racine carrée et cubique dans $\mathbb{R}_+$
\end{itemize}}
& Connaître le vocabulaire lié aux six opérations~:\newline
\textit{Généralités}~: ordre croissant et décroissant, nombre premier, distance à zéro.\newline
\textit{Addition et soustraction}~: terme, somme, différence, opposé.\newline
\textit{Multiplication}~: facteur, produit, multiples, ppmc.\newline
\textit{Division}~: diviseur, dividende, pgdc, numérateur, dénominateur, fraction irréductible, inverse.\newline
\textit{Puissance}~: base, exposant.\newline
\textit{Racine}~: racine carrée et cubique. & & & \\
\cline{2-5}
& Amplifier et simplifier une fraction & & & \\
\cline{2-5}
& Rendre une fraction irréductible & & & \\
\cline{2-5}
& Transformer une fraction en écriture décimale et une écriture décimale en fraction & & & \\
\cline{2-5}
& Classer des nombres réels par ordre croissant ou décroissant en les plaçant sur une droite numérique ou en utilisant les signes $<$, $>$, $\leq$, $\geq$, $=$ & & & \\
\cline{2-5}
& Écrire et lire des nombres écrits sous forme de puissance et de racine & & & \\
\cline{2-5}
& Écrire un nombre rationnel en notation scientifique et vice-versa & & & \\
\cline{2-5}
& Connaître la racine carrée d'un carré parfait ($\leq 144$) et la racine cubique d'un cube parfait ($-216\leq x\leq 216$) & & & \\
\hline
 
\multirow{2}{=}{Discerner les ensembles de nombres, découvrir quelques nombres irrationnels}
& Donner le nom des ensembles de nombres et connaître leur notation~: $\mathbb{N}$, $\mathbb{Z}$, $\mathbb{Q}$, $\mathbb{R}$ & & & \\
\cline{2-5}
& Classer des nombres dans ces différents ensembles & & & \\
\hline
\end{tabularx}
\def\arraystretch{1}
 
%---------------
\newpage
\def\arraystretch{1.6}
\begin{tabularx}{\textwidth}{|>{\raggedright\arraybackslash}p{5.2cm}|>{\raggedright\arraybackslash}X|>{\centering\arraybackslash}p{0.8cm}|>{\centering\arraybackslash}p{0.8cm}|>{\centering\arraybackslash}p{0.8cm}|}
\hline
\rowcolor{themeNO!30}\multicolumn{2}{|l|}{Objectifs} &\multicolumn{3}{c|}{Auto-évaluation} \\
\rowcolor{themeNO!30}\multicolumn{2}{|l|}{} & \Sadey[1.5]& \Neutrey[1.5] & \Smiley[1.5]  \\\hline
 
\multirow{7}{=}{Utiliser des procédures de calcul réfléchi ou de calcul mental pour obtenir un résultat exact ou une estimation, avec des nombres rationnels sous forme décimale ($+$, $-$, $\cdot$, $:$) et sous forme fractionnaire ($+$, $-$, $\cdot$, $:$).
\newline\newline
Utiliser des algorithmes pour effectuer des calculs de façon efficace avec des nombres rationnels sous forme fractionnaire ($+$, $-$, $\cdot$, $:$)}
& Connaître et utiliser les propriétés des six opérations et les priorités des opérations pour effectuer mentalement des chaînes opératoires & & & \\
\cline{2-5}
& Connaître et savoir utiliser en situation les règles d'addition, de soustraction, de multiplication et de division de fractions & & & \\
\cline{2-5}
& Simplifier des fractions quand c'est nécessaire avant d'effectuer les opérations & & & \\
\cline{2-5}
& Déterminer la fraction d'un tout que représente une partie de ce tout & & & \\
\cline{2-5}
& Calculer la fraction ou le pourcentage d'un nombre & & & \\
\cline{2-5}
& Retrouver la valeur du tout lorsqu'on connaît celle d'une fraction de ce tout & & & \\
\cline{2-5}
& Élever une fraction à une puissance & & & \\
\hline
 
\multirow{2}{=}{Explorer des situations aléatoires}
& Connaître le vocabulaire lié aux situations aléatoires~: événement, résultats équiprobables, possibles et favorables & & & \\
\cline{2-5}
& Calculer les chances d'apparition d'un événement~:\newline
en déterminant le nombre de résultats possibles et le nombre de résultats favorables à l'aide d'un tableau ou d'un diagramme en arbre~;\newline
en calculant ensuite le rapport entre les résultats favorables et les résultats possibles & & & \\
\hline
 
\multirow{2}{=}{Prendre en compte l'ordre dans lequel la calculatrice effectue les opérations}
& Choisir les calculs que l'on doit impérativement faire à l'aide de la calculatrice et effectuer mentalement les autres & & & \\
\cline{2-5}
& Donner le résultat le plus précis possible d'un calcul en utilisant à bon escient les diverses fonctions de la calculatrice & & & \\
\hline
\end{tabularx}
\def\arraystretch{1}

%------------------------------------------------
\newpage
\header{Nombres réels - Feuille de route}
 
\newcommand{\rref}[1]{\textcolor{themeColor}{#1}}
 
% Case à cocher
\newcommand{\case}{\raisebox{-0.1ex}{\Large$\square$}}
 
% Marqueurs obligatoire / facultatif (forme, pas seulement couleur)
\newcommand{\oblig}{\tikz[baseline=-0.5ex]{\fill[themeColor] (0,0) circle (0.09);}}
\newcommand{\faclt}{\tikz[baseline=-0.5ex]{\draw[gray!70,line width=0.6pt] (0,0) circle (0.09);}}
 
 \vfill
\begin{tikzpicture}[
circ/.style={circle, fill=themeColor, text=white, font=\bfseries\Large, minimum size=11mm, inner sep=0pt},
 box/.style={rectangle, draw=themeColor!35!white, rounded corners=2pt, fill=themeColor!8!white, inner sep=3mm, text width=14.7cm, align=left, anchor=north west}]
 
% ---------- Étape 1 ----------
\node[circ, anchor=north] (c1) at (0,0) {1};
\node[box] (b1) at ($(c1.north)+(1.2cm,0)$) {%
\textbf{\large Addition et soustraction de fractions}\\[4pt]
{\small\begin{tabularx}{10cm}{cXp{3cm}p{1cm}}
\case & \rref{NO 103} & Fiche, p.\,36 & \oblig \\
\case & \rref{NO 117} & Fiche, p.\,42 & \oblig \\
\case & \rref{NO 122} & Fiche, p.\,42 & \oblig \\
\case & \rref{NO 118} & Livre, p.\,53 & \oblig \\
\case & \rref{NO 119} & Livre, p.\,53 & \oblig \\
\case & NO 120 & Livre, p.\,53 & \faclt \\
\case & NO 121 & Livre, p.\,54 & \faclt \\
\end{tabularx}}
};
 
% ---------- Étape 2 ----------
\coordinate (y2) at ($(b1.south west)+(0,-3mm)$);
\node[circ, anchor=north] (c2) at (c1 |- y2) {2};
\node[box] (b2) at (y2) {%
\textbf{\large Multiplication de fractions}\\[4pt]
{\small\begin{tabularx}{10cm}{cXp{3cm}p{1cm}}
\case & \rref{NO 156} & Fiche, p.\,52 & \oblig \\
\case & \rref{NO 160} & Fiche, p.\,53 & \oblig \\
\case & \rref{NO 162} & Livre, p.\,60 & \oblig \\
\case & NO 163 & Livre, p.\,60 & \faclt \\
\end{tabularx}}
};
 
% ---------- Étape 3 ----------
\coordinate (y3) at ($(b2.south west)+(0,-3mm)$);
\node[circ, anchor=north] (c3) at (c1 |- y3) {3};
\node[box] (b3) at (y3) {%
\textbf{\large Division de fractions}\\[4pt]
{\small\begin{tabularx}{10cm}{cXp{3cm}p{1cm}}
\case & \rref{NO 170} & Fiche, p.\,54 & \oblig \\
\case & \rref{NO 171} & Fiche, p.\,54 & \oblig \\
\case & \rref{NO 172} & Livre, p.\,61 & \oblig \\
\case & NO 173 & Livre, p.\,61 & \faclt \\
\case & NO 174 & Livre, p.\,61 & \faclt \\
\end{tabularx}}
};
 
% ---------- Étape 4 ----------
\coordinate (y4) at ($(b3.south west)+(0,-3mm)$);
\node[circ, anchor=north] (c4) at (c1 |- y4) {4};
\node[box] (b4) at (y4) {%
\textbf{\large Notation scientifique}\\[4pt]
{\small\begin{tabularx}{10cm}{cXp{3cm}p{1cm}}
\case & \rref{NO 210} & Fiche, p.\,63 & \oblig \\
\case & \rref{NO 217} & Livre, p.\,69 & \oblig \\
\case & \rref{NO 212} & Livre, p.\,69 & \oblig \\
\case & \rref{NO 214} & Fiche, p.\,64 & \oblig \\
\case & \rref{NO 215} & Fiche, p.\,65 & \oblig \\
\case & NO 220 & Livre, p.\,70 & \faclt \\
\case & NO 221 & Livre, p.\,70 & \faclt \\
\case & NO 222 & Livre, p.\,70 & \faclt \\
\end{tabularx}}
};
 
% ---------- Étape 5 ----------
\coordinate (y5) at ($(b4.south west)+(0,-3mm)$);
\node[circ, anchor=north] (c5) at (c1 |- y5) {5};
\node[box] (b5) at (y5) {%
\textbf{\large Situations aléatoires}\\[4pt]
{\small\begin{tabularx}{10cm}{cXp{3cm}p{1cm}}
\case & \rref{NO 227} & Livre, p.\,72 & \oblig \\
\case & \rref{NO 228} & Livre, p.\,72 & \oblig \\
\case & \rref{NO 229} & Livre, p.\,72 & \oblig \\
\case & NO 239 & Livre, p.\,76 & \faclt \\
\case & NO 240 & Livre, p.\,76 & \faclt \\
\end{tabularx}}
};
 
% ---------- Ligne de parcours (derrière les cercles) ----------
\begin{pgfonlayer}{background}
\draw[themeColor, line width=2pt] (c1) -- (c2) -- (c3) -- (c4) -- (c5);
\end{pgfonlayer}
 
\end{tikzpicture}
 
 \vfill
\begin{center}
\small \case~à faire \qquad \oblig~\rref{obligatoire} \qquad \faclt~facultatif (consolidation / pour aller plus loin)
\end{center}
\vspace{1em}